% !TeX root = ../../Book2.tex
\chapter{经典逼近论}
\label{b2:ch:approximation}
\BookChapterNumber{b2:ch:approximation}

\begin{chapterabstract}
用多项式逼近连续函数，既要判断误差能否趋于零，也要研究给定次数下能达到多小的误差。本章先以 Bernstein 多项式给出可计算的一致逼近，再用 Chebyshev 交错定理刻画实区间上的最佳多项式。周期卷积的正核构造导出 Jackson 直接定理；三角多项式的导数估计则反过来从逼近速度辨认函数的光滑性。最后讨论高阶差分、整数光滑阶的限制和区间端点的特殊尺度，说明几种看似相近的逼近结论为何不能混用。
\end{chapterabstract}

\begin{learninggoals}
\begin{itemize}
\item\label{b2:item:approximation-goal-modulus}用连续模统一控制紧集上连续函数的局部变化。
\item\label{b2:item:approximation-goal-bernstein}掌握一元与多元 Bernstein 多项式的构造，以非负权重和二阶矩估计证明一致逼近。
\item\label{b2:item:approximation-goal-planned}证明实区间最佳多项式的存在唯一性，用交错极值检验最佳性，并区别一致范数与平方范数的优化目标。
\item\label{b2:item:approximation-goal-jackson}从正核的次数、归一化和一阶矩推出 Jackson 估计，解释构造性逼近与最佳逼近速度的区别。
\item\label{b2:item:approximation-goal-inverse}用 Bernstein 导数估计和二进尺度分解证明周期 Hölder 刻画，辨认整数阶与区间端点的限制。
\end{itemize}
\end{learninggoals}

用有限个数描述一个连续函数，最直接的办法是在一些点上取样。怎样把这些取样值组合成一个多项式，使它在取样点之间也接近原函数？Bernstein 构造把附近的取样值赋予较大权重，提供了一种直观的做法。得到近似之后，还可以继续追问：同样的次数，是否能把最大误差降得更低？这一问题把注意力引向误差最大的那些点，最终产生 Chebyshev 的交错判据。

\ProseParagraphBreak

增加多项式的次数，通常会使近似更精确；但误差究竟下降多快，还取决于原函数怎样变化。平滑的曲线是否比带尖点的曲线更容易逼近？如果已经知道误差下降很快，又能从中读出函数的哪些性质？Jackson 定理与逆定理分别回答这两个方向的问题。章末再考察高阶变化，以及闭区间端点附近与内部不同的逼近尺度。

\ProseParagraphBreak

Elias M. Stein 与 Rami Shakarchi 的《Fourier Analysis: An Introduction》\cite{Stein2003Fourier}可用于回顾前章的周期分析背景；Antoni Zygmund 的《Trigonometric Series》\cite{Zygmund2002Trigonometric}可作为继续研究三角级数与函数光滑性的进阶读物。若想了解区间之外的连续函数逼近，可另读第一卷第29.2节的一般 Stone--Weierstrass 定理。

% !TeX root = ../../Book2.tex
\section{逼近空间、误差与连续模}
\label{b2:sec:approximation-spaces-errors}

\subsection{逼近空间}
\label{b2:subsec:approximation-spaces}

待撰写。

\subsection{连续模与逼近误差}
\label{b2:subsec:020901}

Bernstein 构造在有限个点上取样，再以非负权重求平均。估计取样误差时，需要同时控制所有相距较近的点上的函数值之差。

\begin{definition}
\label{b2:def:approximation-modulus-continuity}
设 \(d\) 为正整数，\(K\subset\mathbb R^d\) 为非空紧集，\(\mathbb F=\mathbb R\) 或 \(\mathbb C\)，且 \(f\in C(K;\mathbb F)\)。对 \(\delta\geq0\)，定义
\[
 \omega_f(\delta)
 \coloneqq\sup_{\substack{x,y\in K\\|x-y|\leq\delta}}
             |f(x)-f(y)|.
\]
这里 \(|x-y|\) 为欧氏距离。函数 \(\omega_f\) 称为 \(f\) 的\term[lianxu-mo]{连续模}[modulus of continuity]。
\end{definition}
\symidx{\(\omega_f\)：函数的连续模}

连续模单调不减，且 \(0\leq\omega_f(\delta)\leq2\lVert f\rVert_\infty\)。紧集上的一致连续性给出 \(\omega_f(\delta)\to0\)（\(\delta\downarrow0\)）；这一步把逐对点的连续性变为统一的误差控制。

\subsection{最佳逼近误差}
\label{b2:subsec:approximation-best-error}

待撰写。

% !TeX root = ../../Book2.tex
\section{Bernstein 构造性逼近}
\label{b2:sec:approximation-bernstein}

要实际计算一个逼近多项式，可以先从等距网格取得函数值
\(f(0),f(1/n),\ldots,f(1)\)，把它们组合成一个明确的多项式。
构造的关键是使各个权重非负、总和为一，而且对固定的 \(x\)，主要权重落在
\(k/n\) 接近 \(x\) 的取样点上。这样，连续性便能把取样点的接近转化为函数值的接近。

\ProseParagraphBreak
这一办法不要求多项式经过所有取样点。插值规定每个节点上的值都完全正确，
这里则允许在节点上有误差，以换取容易控制的加权平均。

\subsection{一元 Bernstein 多项式}
\label{b2:subsec:approximation-bernstein-univariate}

下面的\term[bernstein-bijin-dingli]{Bernstein 逼近定理}[Bernstein approximation theorem]用二项式权重构造连续函数的多项式一致逼近。

\begin{theorem}[Bernstein 逼近：一元情形]
\label{b2:thm:approximation-bernstein-univariate}
设 \(\mathbb F=\mathbb R\) 或 \(\mathbb C\)，且 \(f\in C([0,1];\mathbb F)\)。对每个正整数 \(n\)，定义
\[
 b_{n,k}(x)\coloneqq\binom nk x^k(1-x)^{n-k},
 \quad
 B_nf(x)\coloneqq\sum_{k=0}^nf(k/n)\cdot b_{n,k}(x),
 \quad 0\leq x\leq1.
\]
这里零次幂按 \(1\) 计算。\(B_nf\) 是次数不超过 \(n\) 的\term[bernstein-duoxiangshi]{Bernstein 多项式}[Bernstein polynomial]。以
\[
 \omega_f(\delta)\coloneqq
 \sup\{|f(s)-f(t)|:s,t\in[0,1],\ |s-t|\leq\delta\}
 \quad(\delta>0)
\]
表示 \(f\) 关于通常距离的连续模，则对 \(0<\delta\leq1\)，
\[
 \lVert B_nf-f\rVert_\infty
 \leq\omega_f(\delta)
      +\frac{\lVert f\rVert_\infty}{2n\delta^2}.
\]
因此 \(B_nf\to f\) 一致成立。此外，若 \(a<b\) 为实数，\(h\in C([a,b];\mathbb F)\)，则存在次数不超过 \(n\) 的多项式 \(q_n\)，使 \(q_n\to h\) 在 \([a,b]\) 上一致成立。
\end{theorem}
\begin{proof}
这些权重非负。对固定的 \(x\in[0,1]\)，引入辅助变量 \(z\)，由二项式公式，
\[
 \sum_{k=0}^n b_{n,k}(x)z^k=(1-x+xz)^n.
\]
分别令 \(z=1\)，以及对 \(z\) 求一、二阶导数后令 \(z=1\)，得到
\[
\begin{aligned}
 &\sum_{k=0}^n b_{n,k}(x)=1,\quad
   \sum_{k=0}^n k\cdot b_{n,k}(x)=nx,\\
 &\sum_{k=0}^n k(k-1)\cdot b_{n,k}(x)=n(n-1)x^2.
\end{aligned}
\]
当 \(n=1\) 时，第三个等式两边均为零。这些都是关于 \(x\) 的多项式恒等式，
所以端点 \(x=0,1\) 同样适用。利用 \(k^2=k(k-1)+k\) 展开平方，可得
\begin{equation}
 \label{b2:eq:approximation-bernstein-variance}
 \sum_{k=0}^n(k/n-x)^2\cdot b_{n,k}(x)
 =\frac{x(1-x)}n\leq\frac1{4n}.
\end{equation}
这说明按这些权重平均，取样点与 \(x\) 的平方距离很小。
固定 \(0<\delta\leq1\)，将求和分为 \(|k/n-x|\leq\delta\) 与其补集。邻近部分的函数值之差至多为 \(\omega_f(\delta)\)，而远处权重满足
\[
 \sum_{|k/n-x|>\delta}b_{n,k}(x)
 \leq\frac1{\delta^2}
       \sum_{k=0}^n(k/n-x)^2\cdot b_{n,k}(x)
 \leq\frac1{4n\delta^2}.
\]
由于权重之和为一，
\[
\begin{aligned}
 |B_nf(x)-f(x)|
 &\leq\sum_{k=0}^n|f(k/n)-f(x)|\cdot b_{n,k}(x)\\
 &\leq\omega_f(\delta)
       +2\lVert f\rVert_\infty\cdot\frac1{4n\delta^2}.
\end{aligned}
\]
对 \(x\) 取上确界即得误差估计。取 \(\delta=n^{-1/4}\)，右端成为
\[
 \omega_f(n^{-1/4})+\frac{\lVert f\rVert_\infty}{2\sqrt n},
\]
由一致连续性趋于零。

对一般区间上的 \(h\)，令 \(H(x)=h(a+(b-a)x)\)，其中 \(0\leq x\leq1\)，再设
\[
 q_n(t)\coloneqq B_nH\!\left(\frac{t-a}{b-a}\right),
 \quad a\leq t\leq b.
\]
仿射换元保持次数不超过 \(n\)，而且
\(\lVert q_n-h\rVert_{C[a,b]}=\lVert B_nH-H\rVert_{C[0,1]}\to0\)，得到所需逼近。
\end{proof}

三个权重恒等式还给出可以直接核验的例子。若 \(\iota(x)=x\)，则
\[
 B_n1=1,\quad B_n\iota=\iota,\quad
 B_n(\iota^2)(x)=x^2+\frac{x(1-x)}n.
\]
因此常数和仿射函数被完全保持；对 \(x^2\)，误差在两端为零，在中点达到
\(1/(4n)\)。这也说明 \(B_n\) 一般不是投影到次数不超过 \(n\) 的多项式空间：
即使 \(n\geq2\)，它仍会改变已在该空间中的 \(x^2\)。对所有连续 \(f\)，
端点处只有一个权重非零，所以 \(B_nf(0)=f(0)\)、\(B_nf(1)=f(1)\)；
内部取样点通常没有这种精确相等。

\subsection{连续模与有效取样尺度}
\label{b2:subsec:approximation-bernstein-rate}

前面的远近拆分便于证明收敛，但取 \(\delta=n^{-1/4}\) 并没有反映构造本身的尺度。
由\eqref{b2:eq:approximation-bernstein-variance}，平方距离的平均是 \(O(n^{-1})\)，
所以距离的典型大小应为 \(n^{-1/2}\)。要把这个信息用于误差，
可以将较远的两点之间的变化拆成若干段短距离上的变化。

\begin{lemma}[区间连续模的分段估计]
\label{b2:lem:approximation-modulus-subdivision}
设 \(a<b\) 为实数，\(\mathbb F=\mathbb R\) 或 \(\mathbb C\)，
\(f\in C([a,b];\mathbb F)\)，并令
\(\omega_f(\delta)=\sup_{s,t\in[a,b],\,|s-t|\leq\delta}|f(s)-f(t)|\)。
则对每个 \(\delta>0\) 和 \(s,t\in[a,b]\)，
\[
 |f(s)-f(t)|\leq\left(1+\frac{|s-t|}{\delta}\right)\omega_f(\delta).
\]
\end{lemma}
\begin{proof}
当 \(s=t\) 时左端为零。否则取
\(m=\lceil|s-t|/\delta\rceil\)，把从 \(s\) 到 \(t\) 的线段分成 \(m\) 个等长小段。
每段长不超过 \(\delta\)，所有分点仍在 \([a,b]\) 中。
沿分点逐次使用三角不等式，得到
\(|f(s)-f(t)|\leq m\omega_f(\delta)\)。最后使用
\(m\leq1+|s-t|/\delta\) 即可。
\end{proof}

\begin{proposition}[Bernstein 多项式的连续模估计]
\label{b2:prop:approximation-bernstein-modulus}
设 \(\mathbb F=\mathbb R\) 或 \(\mathbb C\)，\(f\in C([0,1];\mathbb F)\)，
\(n\geq1\) 为整数，且
\[
 B_nf(x)=\sum_{k=0}^nf(k/n)\binom nk x^k(1-x)^{n-k},
\]
端点的零次幂按 \(1\) 计算。以 \(\omega_f\) 表示关于通常距离的连续模。
对每个 \(\delta>0\) 和 \(x\in[0,1]\)，有
\[
 |B_nf(x)-f(x)|
 \leq\left(1+\frac{\sqrt{x(1-x)}}{\sqrt n\,\delta}\right)\omega_f(\delta).
\]
特别地，
\[
 \lVert B_nf-f\rVert_\infty\leq\frac32\omega_f(n^{-1/2}).
\]
若还存在 \(0<\alpha\leq1\) 和 \(L\geq0\)，使
\(|f(s)-f(t)|\leq L|s-t|^\alpha\) 对所有 \(s,t\in[0,1]\) 成立，则
\[
 |B_nf(x)-f(x)|\leq L\left(\frac{x(1-x)}n\right)^{\alpha/2},\quad
 \lVert B_nf-f\rVert_\infty\leq\frac{L}{2^\alpha n^{\alpha/2}}.
\]
\end{proposition}
\begin{proof}
先将\cref{b2:lem:approximation-modulus-subdivision}应用于 \(s=k/n,t=x\)，再按权重求和：
\[
 |B_nf(x)-f(x)|
 \leq\omega_f(\delta)\left(1+\frac1\delta
       \sum_{k=0}^n|k/n-x|\cdot b_{n,k}(x)\right).
\]
由 Cauchy--Schwarz 不等式、权重和为一及\eqref{b2:eq:approximation-bernstein-variance}，
\[
 \begin{aligned}
 \sum_{k=0}^n|k/n-x|\cdot b_{n,k}(x)
 &\leq\left(\sum_{k=0}^n(k/n-x)^2\cdot b_{n,k}(x)\right)^{1/2}
       \left(\sum_{k=0}^nb_{n,k}(x)\right)^{1/2}\\
 &=\sqrt{\frac{x(1-x)}n}.
 \end{aligned}
\]
代回即得第一式。令 \(\delta=n^{-1/2}\)，并使用 \(x(1-x)\leq1/4\)，得到一致范数的估计。

在 Hölder 条件下，直接由有限和的 Hölder 不等式，以共轭指数
\(2/\alpha\) 与 \(2/(2-\alpha)\) 估计，得到
\[
 \begin{aligned}
 |B_nf(x)-f(x)|
 &\leq L\sum_{k=0}^n|k/n-x|^\alpha\cdot b_{n,k}(x)\\
 &\leq L\left(\sum_{k=0}^n(k/n-x)^2\cdot b_{n,k}(x)\right)^{\alpha/2}
       \left(\sum_{k=0}^nb_{n,k}(x)\right)^{1-\alpha/2}.
 \end{aligned}
\]
再用权重恒等式即可。
\end{proof}

这个估计把近似的精度与函数的连续性联系起来：Hölder 指数为 \(\alpha\) 时，
它给出 \(n^{-\alpha/2}\) 的误差上界。它只是这一具体构造的保证，
不是最佳逼近误差的下界，也不能据此断言所有函数的实际误差都恰好有这个量级。
前面的 \(x^2\) 就有更快的 \(n^{-1}\) 误差。\secref{b2:sec:approximation-jackson}将用更集中的核，
把一般 Hölder 函数的逼近上界改进到 \(n^{-\alpha}\)。

\subsection{正性、稳定性与形状保持}
\label{b2:subsec:approximation-bernstein-shape}

非负权重除了简化误差估计，还使构造具有直接的稳定性。
如果每个取样值都有至多 \(\varepsilon\) 的误差，加权平均后的误差仍不超过
\(\varepsilon\)，不会因节点数增加而放大。对于实值函数，
这些权重还保留单调性与凸性。

\begin{proposition}[Bernstein 算子的稳定性与保形性]
\label{b2:prop:approximation-bernstein-shape}
设 \(\mathbb F=\mathbb R\) 或 \(\mathbb C\)，\(n\geq1\) 为整数，并对
\(f\in C([0,1];\mathbb F)\) 定义
\[
 B_nf(x)=\sum_{k=0}^nf(k/n)\binom nk x^k(1-x)^{n-k},\quad 0\leq x\leq1,
\]
端点的零次幂按 \(1\) 计算。则 \(B_n\) 是线性映射，且对所有
\(f,g\in C([0,1];\mathbb F)\)，
\[
 \lVert B_nf-B_ng\rVert_\infty\leq\lVert f-g\rVert_\infty.
\]
当 \(\mathbb F=\mathbb R\) 时，\(f\geq0\) 蕴含 \(B_nf\geq0\)；
若 \(f\) 单调不减，则 \(B_nf\) 单调不减；若 \(f\) 为凸函数，则 \(B_nf\) 为凸函数，
并且 \(B_nf(x)\geq f(x)\) 对全部 \(x\in[0,1]\) 成立。
\end{proposition}
\begin{proof}
线性由有限求和直接给出。由非负权重之和为一，
\[
 |B_nf(x)-B_ng(x)|
 \leq\sum_{k=0}^n|f(k/n)-g(k/n)|\cdot b_{n,k}(x)
 \leq\lVert f-g\rVert_\infty.
\]
对 \(x\) 取上确界即得稳定性；正性同样由各项非负得到。

为考察形状，暂将越界的 \(b_{m,k}\) 定义为零，并令 \(b_{0,0}=1\)。
对权重求导并利用二项式系数的关系，有
\[
 b'_{n,k}=n(b_{n-1,k-1}-b_{n-1,k}).
\]
两部分分别换指标，端点的越界项为零，便得到
\[
 (B_nf)'(x)=n\sum_{k=0}^{n-1}
       \bigl(f((k+1)/n)-f(k/n)\bigr)b_{n-1,k}(x).
\]
当 \(f\) 单调不减时，每个差均非负，所以导数非负。
当 \(n\geq2\) 时，记
\(d_k=f((k+2)/n)-2f((k+1)/n)+f(k/n)\)。再次求导，得到
\[
 (B_nf)''(x)=n(n-1)\sum_{k=0}^{n-2}d_k\cdot b_{n-2,k}(x).
\]
对凸函数，中点处的值不超过两端值的平均，所以二阶差 \(d_k\) 非负，
从而二阶导数非负。\(n=1\) 时 \(B_1f\) 为仿射函数，也为凸函数。
最后，由有限凸组合的凸性不等式及 \(\sum_k(k/n)b_{n,k}(x)=x\)，有
\[
 f(x)=f\!\left(\sum_{k=0}^n(k/n)\cdot b_{n,k}(x)\right)
 \leq\sum_{k=0}^nf(k/n)\cdot b_{n,k}(x)=B_nf(x).
\]
有限凸组合的这一不等式由凸函数定义中的两点不等式逐次合并得到，
零权重项直接略去即可。
\end{proof}

取 \(f(x)=|x-1/2|\)，原函数在中点不可微，而 \(B_nf\) 是处处光滑的凸多项式，
并始终位于 \(f\) 的上方。这里不需要对 \(f\) 求导：多项式的导数来自离散取样值的差。
因此连续函数的某些几何性质，可以在它尚无导数时就传给光滑近似。

\subsection{多元 Bernstein 多项式}
\label{b2:subsec:approximation-bernstein-multivariate}

在立方体上，可沿每个坐标分别使用同样的二项式权重，再取乘积。
所得网格共有 \((n+1)^d\) 个取样点，每个坐标的集中程度仍由一元计算控制。
这里要区分每个变量的次数与总次数：乘积可能同时含有各坐标的 \(n\) 次幂，
所以总次数上界是 \(dn\)，而不是 \(n\)。

\begin{theorem}[Bernstein 逼近：多元情形]
\label{b2:thm:approximation-bernstein-multivariate}
设 \(d\) 为正整数，\(\mathbb F=\mathbb R\) 或 \(\mathbb C\)，且 \(f\in C([0,1]^d;\mathbb F)\)。对正整数 \(n\)、\(\mathbf k=(k_1,\ldots,k_d)\in\{0,\ldots,n\}^d\) 与 \(x\in[0,1]^d\)，定义
\[
\begin{aligned}
 w_{n,\mathbf k}(x)
 &\coloneqq\prod_{j=1}^d\binom n{k_j}x_j^{k_j}(1-x_j)^{n-k_j},\\
 B_nf(x)
 &\coloneqq\sum_{\mathbf k\in\{0,\ldots,n\}^d}
       f(\mathbf k/n)\cdot w_{n,\mathbf k}(x).
\end{aligned}
\]
零次幂按 \(1\) 计算。\(B_nf\) 在每个变量上的次数均不超过 \(n\)，总次数不超过 \(dn\)。以
\[
 \omega_f(\delta)\coloneqq
 \sup\{|f(s)-f(t)|:s,t\in[0,1]^d,\ |s-t|\leq\delta\}
\]
表示 \(f\) 关于欧氏距离的连续模，则对每个 \(\delta>0\)，
\[
 \lVert B_nf-f\rVert_\infty
 \leq\omega_f(\delta)
       +\frac{d\lVert f\rVert_\infty}{2n\delta^2}.
\]
特别地，\(B_nf\to f\) 在 \([0,1]^d\) 上一致成立；还有较细的估计
\[
 \lVert B_nf-f\rVert_\infty\leq\frac32\omega_f(\sqrt{d/n}).
\]
\end{theorem}
\begin{proof}
乘积权重非负。对每个坐标依次求和，并使用一元二项式恒等式，得到
\[
 \sum_{\mathbf k}w_{n,\mathbf k}(x)=1,
 \quad
 \sum_{\mathbf k}\frac{k_j}{n}\cdot w_{n,\mathbf k}(x)=x_j,
 \quad 1\leq j\leq d.
\]
同样计算每个坐标的二阶矩，
\[
 \sum_{\mathbf k}(k_j/n-x_j)^2\cdot w_{n,\mathbf k}(x)
 =\frac{x_j(1-x_j)}n.
\]
把这些等式相加，便有
\[
 \sum_{\mathbf k}|\mathbf k/n-x|^2\cdot w_{n,\mathbf k}(x)
 =\sum_{j=1}^d\frac{x_j(1-x_j)}n
 \leq\frac d{4n}.
\]
所以对任意 \(\delta>0\)，远处取样点的总权重满足
\[
 \sum_{|\mathbf k/n-x|>\delta}w_{n,\mathbf k}(x)
 \leq\frac d{4n\delta^2}.
\]
在 \(|\mathbf k/n-x|\leq\delta\) 的部分，用连续模控制函数值之差；在补集上，用 \(2\lVert f\rVert_\infty\) 控制。由此得到
\[
\begin{aligned}
 |B_nf(x)-f(x)|
 &\leq\sum_{\mathbf k}|f(\mathbf k/n)-f(x)|\cdot w_{n,\mathbf k}(x)\\
 &\leq\omega_f(\delta)
       +2\lVert f\rVert_\infty\cdot\frac d{4n\delta^2}.
\end{aligned}
\]
取上确界即得所述估计。先由一致连续性取足够小的 \(\delta\)，再令 \(n\) 足够大，就得到一致收敛。

较细的估计使用与一元情形相同的分段办法。立方体是凸集，
连接其中两点的线段仍留在立方体内，故按长度至多 \(\delta\) 分段可以得到
\[
 |f(s)-f(t)|\leq\left(1+\frac{|s-t|}{\delta}\right)\omega_f(\delta).
\]
代入加权和，再用 Cauchy--Schwarz 不等式及上述平方距离估计，得到
\[
 \begin{aligned}
 |B_nf(x)-f(x)|
 &\leq\left(1+\frac1\delta
 \sqrt{\sum_{\mathbf k}|\mathbf k/n-x|^2\cdot w_{n,\mathbf k}(x)}\right)
 \omega_f(\delta)\\
 &\leq\left(1+\frac{\sqrt d}{2\sqrt n\,\delta}\right)\omega_f(\delta).
 \end{aligned}
\]
取 \(\delta=\sqrt{d/n}\) 并对 \(x\) 取上确界，便得第二个估计。
\end{proof}

对于任意非退化闭矩形 \(Q=\prod_{j=1}^d[a_j,b_j]\)，逐坐标作仿射换元
\(x_j=(t_j-a_j)/(b_j-a_j)\)，便把任意 \(C(Q)\) 函数转成单位立方体上的连续函数。
把所得 Bernstein 多项式换回原变量，仍是多项式，且一致误差完全相同。
所以多项式在 \(C(Q)\) 中一致稠密。这个推论正是后文证明某些正交展开完备性时所需的具体工具。

\ProseParagraphBreak
若函数最初只定义在球面 \(S^{d-1}\) 上，则还不能直接对它取立方体网格上的值。
对 \(d\geq2\) 和 \(f\in C(S^{d-1};\mathbb F)\)，其中 \(\mathbb F=\mathbb R\) 或 \(\mathbb C\)，
可取连续函数 \(\chi:[0,\infty)\to[0,1]\)，使
\(\chi(1)=1\)、\(\chi(r)=0\) 于 \(r\leq1/2\)，并令
\[
 F(x)=\chi(|x|)\cdot f(x/|x|)\quad(x\ne0),\quad F(0)=0.
\]
这个函数在原点附近恒为零，在其余位置由连续函数复合而成，因而在包住球面的立方体上连续，
且限制到球面恰好等于 \(f\)。对 \(F\) 作多项式一致逼近，再限制到球面，误差只会减小。
\secref{b2:sec:B203}的球谐展开将实际使用这条路径：先作连续延拓，
再作多项式逼近，最后处理多项式的调和分解。这样既说明多元定理的用途，
也交代了从矩形到球面所需的额外一步。

% !TeX root = ../../Book2.tex
\section{一致最佳逼近与 Chebyshev 理论}
\label{b2:sec:approximation-chebyshev}
\label{b2:sec:0208}

待撰写。

\subsection{一致范数下的最佳逼近}
\label{b2:subsec:020801}

待撰写。

\subsection{Chebyshev 多项式与极值性质}
\label{b2:subsec:020802}

待撰写。

\subsection{交错定理与唯一性}
\label{b2:subsec:020803}

待撰写。

\subsection{与正交最佳逼近的比较}
\label{b2:subsec:020804}

待撰写。

% !TeX root = ../../Book2.tex
\section{Jackson 定理与直接逼近}
\label{b2:sec:approximation-jackson}
\label{b2:sec:0209}

Bernstein 多项式给出了逐点取样的显式算法。一致收敛以后，还有一个更细的问题：若已知函数在距离 \(\delta\) 内的振幅不超过 \(\omega_f(\delta)\)，次数为 \(n\) 的多项式能否把误差降到 \(\omega_f(1/n)\) 的量级？本节先在圆周上用卷积构造这种逼近，再通过余弦变换处理闭区间。

\ProseParagraphBreak

这里的\term[zhijie-bijin-dingli]{直接逼近定理}[direct approximation theorem]是指由函数的连续性或光滑性推出逼近误差的上界。证明要同时保住两个条件：卷积核的质量足够集中，所得函数的次数也不能超过预算。只让核越来越窄，而不计算它的次数，并不能回答给定 \(n\) 时的问题。

\subsection{Jackson 核与直接定理}
\label{b2:subsec:020902}

回顾\subsecref{b2:subsec:020401}的约定：圆周为 \(\mathbb T=\mathbb R/\mathbb Z\)，积分采用总质量为 \(1\) 的 Lebesgue 测度，而 Fejér 核 \(F_N\) 的次数为 \(N\)。因此，若用正整数 \(m\) 表示集中尺度，应该写
\[
 F_{m-1}(t)=\frac1m
  \left(\frac{\sin(m\pi t)}{\sin(\pi t)}\right)^2,
 \quad F_{m-1}(0)=m.
\]
Fejér 核已经非负且积分为 \(1\)，但它离开原点以后只按 \(|t|^{-2}\) 衰减。要估计连续模，需要控制核的第一矩；这个衰减速度会带来一个对数。把核平方，再恢复总质量为 \(1\)，可以改善尾部而不失去非负性。

\begin{definition}
\label{b2:def:approximation-jackson-kernel}
设 \(m\geq1\) 为整数，\(F_{m-1}\) 是单位周期圆周 \(\mathbb T=\mathbb R/\mathbb Z\) 上的 Fejér 核。记
\[
 A_m\coloneqq\int_{\mathbb T}F_{m-1}(t)^2\,\mathrm dt,
 \quad
 J_m(t)\coloneqq\frac{F_{m-1}(t)^2}{A_m}.
\]
函数 \(J_m\) 称为本节使用的\term[jackson-he]{Jackson 核}[Jackson kernel]。
\end{definition}
\symidx{\(J_m\)：Jackson 核}

\begin{lemma}
\label{b2:lem:approximation-jackson-kernel}
设 \(m\geq1\) 为整数，\(F_{m-1}\) 是 \(\mathbb T=\mathbb R/\mathbb Z\) 上的 Fejér 核，积分取归一化 Lebesgue 测度，\(A_m=\int_{\mathbb T}F_{m-1}^2\)、\(J_m=F_{m-1}^2/A_m\)。则 \(J_m\) 是实值、偶、非负的三角多项式，次数不超过 \(2(m-1)\)，且
\[
 \int_{\mathbb T}J_m(t)\,\mathrm dt=1,
 \quad A_m=\frac{2m^2+1}{3m}.
\]
存在与 \(m\) 无关的常数 \(C>0\)，使
\[
 \begin{aligned}
 J_m(t)&\leq\frac{Cm}{(1+m|t|)^4}
 \quad(-1/2\leq t\leq1/2),\\
 \int_{-1/2}^{1/2}|t|J_m(t)\,\mathrm dt&\leq\frac Cm.
 \end{aligned}
\]
这里 \(|t|\) 就是所选代表到圆周原点的最短弧距。
\end{lemma}
\begin{proof}
Fejér 核的有限展开为
\[
 F_{m-1}(t)=\sum_{|k|<m}\left(1-\frac{|k|}{m}\right)e^{2\pi\mathrm i kt}.
\]
两个次数不超过 \(m-1\) 的三角多项式相乘，频率相加，所以平方的次数不超过 \(2(m-1)\)。实值、偶性和非负性直接保留。利用有限指数函数的正交性，有
\[
 A_m=1+2\sum_{k=1}^{m-1}\left(1-\frac{k}{m}\right)^2
 =1+\frac{2}{m^2}\cdot\frac{(m-1)m(2m-1)}6
 =\frac{2m^2+1}{3m}.
\]
因而 \(2m/3\leq A_m\leq m\)，特别地，归一化分母严格为正。

从有限指数和的平方模表达得到 \(F_{m-1}\leq m\)。另一方面，在 \(0<|t|\leq1/2\) 上，\(|\sin(\pi t)|\geq2|t|\) 且 \(|\sin(m\pi t)|\leq1\)，所以
\[
 0\leq F_{m-1}(t)\leq
 \min\left\{m,\frac1{4mt^2}\right\}.
\]
平方后除以 \(A_m\geq2m/3\)，便有
\[
 J_m(t)\leq C m\min\{1,(m|t|)^{-4}\}
 \leq\frac{C' m}{(1+m|t|)^4}.
\]
最后令 \(u=m|t|\)，把正半轴上的积分上界延长到无穷，得到
\[
 \int_{-1/2}^{1/2}|t|J_m(t)\,\mathrm dt
 \leq\frac{2C'}m\int_0^\infty\frac{u}{(1+u)^4}\,\mathrm du
 \leq\frac{C''}m.
\]
积分有限，因为被积函数在无穷远按 \(u^{-3}\) 衰减。质量为 \(1\) 则来自 \(A_m\) 的定义。
\end{proof}

核的平方把次数增加一倍，也把远处衰减从二次提高到四次；归一化以后，它的高度仍为 \(m\) 的量级，主要质量仍落在宽度约 \(1/m\) 的区间内。第一矩估计量化了这个宽度，下面的误差证明只需要质量与第一矩这两项信息。

\begin{theorem}[Jackson：周期情形]
\label{b2:thm:approximation-jackson-periodic}
设 \(\mathbb F=\mathbb R\) 或 \(\mathbb C\)，\(f\in C(\mathbb T;\mathbb F)\)，其中 \(\mathbb T=\mathbb R/\mathbb Z\) 取归一化 Lebesgue 测度。以 \(\omega_f\) 表示关于最短弧距的连续模，以 \(\mathcal E_n(f)\) 表示次数不超过 \(n\) 的三角多项式的一致最佳逼近误差。存在绝对常数 \(C>0\)，使对每个整数 \(n\geq1\)，都有
\[
 \mathcal E_n(f)\leq C\omega_f(1/n).
\]
具体地，取 \(m=\lfloor n/2\rfloor+1\)，则
\[
 Q_nf\coloneqq f*J_m
\]
是次数不超过 \(n\) 的三角多项式，并满足
\[
 \lVert f-Q_nf\rVert_\infty\leq C\omega_f(1/n),
 \quad \lVert Q_nf\rVert_\infty\leq\lVert f\rVert_\infty.
\]
算子 \(Q_n\) 保持常数；对实值函数还保持非负性。
\end{theorem}
\begin{proof}
把 \(J_m\) 展为有限 Fourier 和，逐项积分即可写出
\[
 (f*J_m)(x)=\sum_{|k|\leq2(m-1)}
     \widehat f(k)\widehat J_m(k)e^{2\pi\mathrm i kx}.
\]
这里不需要 \(f\) 本身的 Fourier 级数收敛，因为被展开的是一个有限多项式。由 \(2(m-1)\leq n\)，所得近似函数符合次数限制。

利用 \(J_m\geq0\) 和 \(\int J_m=1\)，有
\[
 |f(x)-(f*J_m)(x)|
 \leq\int_{-1/2}^{1/2}|f(x)-f(x-t)|J_m(t)\,\mathrm dt.
\]
沿圆周最短弧把长度 \(|t|\) 分成至多 \(1+m|t|\) 段，每段长度不超过 \(1/m\)，由三角不等式得到
\[
 |f(x)-f(x-t)|\leq(1+m|t|)\omega_f(1/m).
\]
代回后，\cref{b2:lem:approximation-jackson-kernel}给出
\[
 \lVert f-f*J_m\rVert_\infty
 \leq\omega_f(1/m)\left(1+m\int_{-1/2}^{1/2}|t|J_m(t)\,\mathrm dt\right)
 \leq C\omega_f(1/m).
\]
由于 \(m\geq n/2\)，再把长度不超过 \(2/n\) 的短弧分为两段，便有 \(\omega_f(1/m)\leq2\omega_f(1/n)\)。这些估计也包括 \(n=1\)：此时 \(m=1\)、\(J_1=1\)，输出就是 \(f\) 的平均值。收缩性、保持常数和保持非负性都由非负核的单位质量直接得到。
\end{proof}

这一结果称为\term[jackson-dingli]{Jackson 定理}[Jackson's theorem]。构造出的 \(Q_nf\) 给出了最佳误差的一个上界；具体函数的最佳误差还可能更小。下面将它与前章的 Fejér 平均比较，看看核的选择怎样影响收敛速度。

\ProseParagraphBreak

与 Fejér 平均的差别可以在一个具体函数上看清。取 \(f(t)=\mathrm d_{\mathbb T}(t,0)\)，它是 Lipschitz 常数为 \(1\) 的连续周期函数。在原点，Fejér 逼近的误差恰为
\[
 (f*F_{m-1})(0)-f(0)
 =\int_{-1/2}^{1/2}|t|F_{m-1}(t)\,\mathrm dt.
\]
将积分分为 \(|t|\leq1/m\) 与其余部分，上述核估计给出上界 \(C\log(m+1)/m\)。这个对数不是估计方法造成的假象：对 \(m\geq4\)，在
\[
 I_k=\left[\frac{k+1/4}{m},\frac{k+3/4}{m}\right],
 \quad 1\leq k\leq\lfloor m/2\rfloor-1,
\]
有 \(\sin^2(m\pi t)\geq1/2\)，而 \(\sin^2(\pi t)\leq\pi^2t^2\)。因而每个 \(I_k\) 对第一矩的贡献至少为 \(c/(m(k+1))\)；对 \(k\) 求和得到 \(c'\log(m+1)/m\) 的下界。相比之下，Jackson 构造对同一函数的误差至多为 \(C/n\)。Fejér 平均很好地解决了收敛问题，若还要得到一般 Lipschitz 函数的最佳次数量级，就需要更集中的核。

\subsection{代数多项式的直接逼近}
\label{b2:subsec:approximation-jackson-algebraic}

闭区间上的函数一般不满足端点值相等，直接作周期延拓可能产生跳跃。余弦变换提供另一条路线：它让圆周先沿区间走一遍，再沿原路返回，因而两端都自动接上。

\begin{theorem}[Jackson：闭区间情形]
\label{b2:thm:approximation-jackson-algebraic}
设 \(a<b\) 为实数，\(\mathbb F=\mathbb R\) 或 \(\mathbb C\)，且 \(f\in C([a,b];\mathbb F)\)。令 \(E_n(f)\) 为次数不超过 \(n\) 的代数多项式在 \([a,b]\) 上的一致最佳逼近误差，\(\omega_f\) 为通常距离下的连续模。存在绝对常数 \(C>0\)，使对所有整数 \(n\geq1\)，
\[
 E_n(f)\leq C\omega_f\left(\frac{b-a}{n}\right).
\]
特别地，若 \(0<\alpha\leq1\)、\(H\geq0\)，且 \(|f(x)-f(y)|\leq H|x-y|^\alpha\) 对所有 \(x,y\in[a,b]\) 成立，则
\[
 E_n(f)\leq CH\left(\frac{b-a}{n}\right)^\alpha.
\]
\end{theorem}
\begin{proof}
设 \(c=(a+b)/2\)、\(r=(b-a)/2\)，并在单位圆周上定义
\[
 g(t)=f(c+r\cos(2\pi t)).
\]
这是连续、偶的周期函数。余弦的导数绝对值不超过 \(1\)，所以沿最短圆弧取代表后，有
\[
 |(c+r\cos(2\pi t))-(c+r\cos(2\pi s))|
 \leq2\pi r\mathrm d_{\mathbb T}(t,s).
\]
从而 \(\omega_g(\delta)\leq\omega_f(\pi(b-a)\delta)\)。这里连续模的参数允许超过区间长度；超过以后，它仅仅等于函数在整个区间上的最大振幅。

对 \(g\) 使用周期 Jackson 构造，得到 \(q_n=g*J_m\)，其中 \(m=\lfloor n/2\rfloor+1\)。核与 \(g\) 都是偶函数，换元 \(t\mapsto-t\) 便知 \(q_n\) 也偶，因此它的有限三角展开只有余弦项：
\[
 q_n(t)=a_0+\sum_{k=1}^na_k\cos(2\pi kt).
\]
余弦递推式
\[
 \cos((k+1)\theta)=2\cos\theta\cos(k\theta)-\cos((k-1)\theta)
\]
表明 \(\cos(k\theta)=T_k(\cos\theta)\)，其中 \(T_0=1\)、\(T_1(u)=u\)、\(T_{k+1}(u)=2uT_k(u)-T_{k-1}(u)\)。因此
\[
 p_n(x)=a_0+\sum_{k=1}^na_kT_k\left(\frac{x-c}{r}\right)
\]
是次数不超过 \(n\) 的代数多项式，而且 \(p_n(c+r\cos(2\pi t))=q_n(t)\)。由于这个余弦参数化覆盖整个 \([a,b]\)，两种一致误差恰好相同。于是
\[
 \lVert f-p_n\rVert_{\infty,[a,b]}
 =\lVert g-q_n\rVert_{\infty,\mathbb T}
 \leq C\omega_g(1/n)
 \leq C\omega_f\left(\frac{\pi(b-a)}n\right).
\]
在区间中把任意一段分为四段，连续模的次可加性给出
\(\omega_f(\pi\delta)\leq4\omega_f(\delta)\)。把 \(4\) 吸收入绝对常数，就得到第一式；代入 Hölder 条件即得第二式。
\end{proof}

例如 \(f(x)=|x|\) 在 \([-1,1]\) 上的连续模满足 \(\omega_f(\delta)\leq\delta\)，所以存在次数为 \(n\) 的代数多项式使一致误差不超过 \(C/n\)，尽管 \(f\) 在原点不可微。这里的构造先把尖点变为 \(|\cos(2\pi t)|\) 的周期尖点，用 Jackson 核平滑，再把偶三角多项式还原为代数多项式。

\ProseParagraphBreak

Bernstein 取样构造与 Jackson 构造的用途有所不同：前者直接使用网格上的函数值，并有清楚的非负权重；后者通过卷积得到更好的通用次数估计。下一节把问题反过来：如果最佳误差本来就衰减得快，函数本身是否必须有相应的连续模？回答这个问题需要先限制有限次数多项式的变化速度。

% !TeX root = ../../Book2.tex
\section{Bernstein 估计与逆定理}
\label{b2:sec:approximation-inverse}

待撰写。

\subsection{三角多项式的 Bernstein 估计}
\label{b2:subsec:020903}

待撰写。

\subsection{逼近速度与光滑性的逆向关系}
\label{b2:subsec:approximation-inverse-smoothness}

待撰写。

% !TeX root = ../../Book2.tex
\OptionalSection{高阶逼近与端点问题}{b2:sec:approximation-higher-order}

一阶连续模只测量函数值的变化。即使函数无限次可微，只要不是常数，它的一阶连续模通常也只有一阶小量；因而仅用它不能分辨二阶、三阶光滑性。上节在 \(0<\alpha<1\) 时把逼近速度与 Hölder 连续性联系起来；光滑阶达到整数时，二者的关系还会出现新的现象。本节先说明高阶差分怎样消去低阶变化，再分别处理光滑指数的整数端点和空间区间的几何端点。这两种“端点”是不同的问题。

\subsection{高阶连续模}
\label{b2:subsec:approximation-higher-moduli}

从二阶开始看最直接。若 \(f\) 在一小段上近似为仿射函数，则 \(f(x+2h)-2f(x+h)+f(x)\) 比 \(f(x+h)-f(x)\) 更小：常数项和一次项恰好抵消。高阶差分把这一消去过程重复进行，而不要求预先知道导数存在。

\begin{definition}\label{b2:def:approximation-higher-modulus}
设 \(\mathbb F=\mathbb R\) 或 \(\mathbb C\)，\(f\in C(\mathbb T;\mathbb F)\)，其中 \(\mathbb T=\mathbb R/\mathbb Z\)，且 \(r\geq1\) 为整数。对 \(h\in\mathbb R\)，定义 \(r\) 阶差分及相应的\term[gaojie-lianxu-mo]{高阶连续模}[higher order modulus of smoothness]为
\[
 \Delta_h^r f(x)=\sum_{j=0}^r(-1)^{r-j}\binom rj f(x-jh),
 \quad
 \omega_r(f,\delta)=\sup_{|h|\leq\delta}\|\Delta_h^r f\|_\infty
 \quad(\delta\geq0).
\]
自变量均按模 \(1\) 理解。特别地，\(\omega_1(f,\delta)=\omega_f(\delta)\) 是圆周上的一阶连续模。
\end{definition}
\symidx{\(\omega_r(f,\delta)\)：周期函数的高阶连续模}

这里允许负步长，且 \(\delta\) 可以大于半个周期；这使伸缩估计无需另外分段。\(\Delta_h^r=(\tau_h-I)^r\)，其中 \(\tau_hf(x)=f(x-h)\)，所以差分是平移的有限组合。连续性与紧性保证上式范数有限，且 \(\omega_r(f,\delta)\to0\) 当 \(\delta\downarrow0\)。

\begin{proposition}\label{b2:prop:approximation-difference-estimates}
设 \(r\geq1\) 为整数，\(\mathbb F=\mathbb R\) 或 \(\mathbb C\)，且 \(f\in C(\mathbb T;\mathbb F)\)。对 \(\delta>0\)、\(\lambda>0\)，有
\[
 \omega_r(f,\lambda\delta)\leq(1+\lambda)^r\omega_r(f,\delta),
 \quad
 \omega_r(f,\delta)\leq2^{r-1}\omega_1(f,\delta).
\]
若进一步 \(f\in C^r(\mathbb T;\mathbb F)\)，则
\[
 \omega_r(f,\delta)\leq\delta^r\|f^{(r)}\|_\infty.
\]
\end{proposition}
\begin{proof}
对任意整数 \(m\geq1\)，平移算子相互交换，并且
\[
 \tau_{mh}-I=(I+\tau_h+\cdots+\tau_{(m-1)h})(\tau_h-I).
\]
取 \(r\) 次幂。每个平移在一致范数下都是等距算子，括号中有 \(m\) 项，故 \(\|\Delta_{mh}^r f\|_\infty\leq m^r\|\Delta_h^r f\|_\infty\)。对 \(|u|\leq\lambda\delta\)，取 \(m=\max\{1,\lceil\lambda\rceil\}\)、\(h=u/m\)，得到第一式。第二式来自 \(\Delta_h^r=(\tau_h-I)^{r-1}\Delta_h\) 和 \(\|\tau_h-I\|\leq2\)。

对 \(C^r\) 函数，反复使用微积分基本定理得到
\[
 \Delta_h^r f(x)=\int_0^h\cdots\int_0^h
 (-1)^rf^{(r)}(x-t_1-\cdots-t_r)\,\mathrm dt_1\cdots\mathrm dt_r.
\]
负的 \(h\) 按有向积分解释。取绝对值后，积分区域的体积是 \(|h|^r\)，于是得到最后的估计。
\end{proof}

高阶差分消去了低阶变化，也改变了误差的估计方式。为了利用这种抵消，需要让核的尾部衰减得更快，再用它平均高阶差分。将 Fejér 核取更高的整数次幂，就能得到所需的衰减。

\begin{theorem}[Jackson：高阶周期形式]\label{b2:thm:approximation-jackson-higher}
对每个整数 \(r\geq1\)，存在只依赖 \(r\) 的常数 \(C_r>0\)，使任意 \(\mathbb F=\mathbb R\) 或 \(\mathbb C\)、\(f\in C(\mathbb T;\mathbb F)\) 及整数 \(n\geq1\) 均满足
\[
 \mathcal E_n(f)\leq C_r\omega_r(f,1/n).
\]
这里 \(\mathcal E_n(f)\) 是用次数不超过 \(n\) 的三角多项式作一致逼近的最佳误差。因此，若 \(f\in C^r(\mathbb T;\mathbb F)\)，则 \(\mathcal E_n(f)\leq C_r n^{-r}\|f^{(r)}\|_\infty\)。
\end{theorem}
\begin{proof}
固定整数 \(s\geq2\)，使 \(2s>r+1\)。沿用\cref{b2:def:02-fejer-kernel}的编号，\(F_{m-1}\) 的次数为 \(m-1\)。对 \(m\geq1\)，定义
\[
 a_{m,s}=\int_{-1/2}^{1/2}F_{m-1}(t)^s\,\mathrm dt,
 \quad K_{m,s}(t)=\frac{F_{m-1}(t)^s}{a_{m,s}}.
\]
这是非负、积分为 \(1\)、次数为 \(s(m-1)\) 的三角多项式。

先核对所需估计。由正弦商公式，在 \(|t|\leq1/2\) 上
\[
 F_{m-1}(t)\leq\min\left\{m,\frac1{4mt^2}\right\}
 \leq\frac{C m}{(1+m|t|)^2}.
\]
在 \(|t|\leq1/(2m)\) 上，用 \(|\sin(m\pi t)|\geq2m|t|\)、\(|\sin(\pi t)|\leq\pi|t|\)，可得 \(F_{m-1}(t)\geq4m/\pi^2\)，零点处按连续性理解。因而下界通过在这段长 \(1/m\) 的区间积分取得，上界通过代换 \(u=mt\) 取得：
\[
 c_s m^{s-1}\leq a_{m,s}\leq C_s m^{s-1},
 \quad K_{m,s}(t)\leq\frac{C_s m}{(1+m|t|)^{2s}}.
\]
这些估计也包含 \(m=1\)。由 \(2s>r+1\)，进一步有
\[
 \int_{-1/2}^{1/2}(1+m|t|)^rK_{m,s}(t)\,\mathrm dt\leq C_r.
\]
这里及以下选定的 \(s\) 只依赖 \(r\)，常数可统一记为 \(C_r\)。

现在取 \(m=\lfloor n/s\rfloor+1\)，使 \(s(m-1)\leq n\) 且 \(n/m\leq s\)，并令
\[
 Q_n(x)=f(x)-(-1)^r\int_{-1/2}^{1/2}\Delta_t^r f(x)K_{m,s}(t)\,\mathrm dt.
\]
为什么这个表达式是三角多项式，而不是仍含未知频率的 \(f\)？差分中 \(j=0\) 的项是 \((-1)^rf(x)\)，积分后恰好与前面的 \(f(x)\) 抵消。因此 \(Q_n\) 是下列 \(r\) 个函数的有限线性组合：
\[
 G_j(x)=\int_{\mathbb T}f(x-jt)K_{m,s}(t)\,\mathrm dt,
 \quad 1\leq j\leq r.
\]
连续函数在紧圆周上有界，故 Fubini 可用于计算其 Fourier 系数。换元 \(y=x-jt\) 给出
\[
 \widehat G_j(k)=\widehat f(k)\widehat K_{m,s}(jk).
\]
当 \(|jk|>s(m-1)\) 时右端为零。取由其有限个非零系数组成的三角多项式，它与 \(G_j\) 的全部 Fourier 系数一致；\cref{b2:cor:02-uniqueness}给出几乎处处相等，而二者连续，故处处相等。这样 \(G_j\) 的次数不超过 \(\lfloor s(m-1)/j\rfloor\)，\(Q_n\) 的次数不超过 \(n\)。实值输入也给出实值 \(Q_n\)。

最后直接从构造估计误差。由连续模的伸缩估计和核的非负性，
\[
 \begin{aligned}
 \|f-Q_n\|_\infty
 &\leq\int_{-1/2}^{1/2}\omega_r(f,|t|)K_{m,s}(t)\,\mathrm dt\\
 &\leq\omega_r(f,1/n)
       \int_{-1/2}^{1/2}(1+n|t|)^rK_{m,s}(t)\,\mathrm dt
 \leq C_r\omega_r(f,1/n).
 \end{aligned}
\]
最后一步使用 \(n/m\leq s\) 及刚建立的矩估计。取最佳误差即得结论，再用\cref{b2:prop:approximation-difference-estimates}得到 \(C^r\) 情形。
\end{proof}

这里 \(Q_n\) 在 \(r\geq2\) 时一般不是正平均；正的是用来估计差分的核。正性负责误差控制，差分的系数负责消去低阶变化，二者承担不同任务。差分估计还适用于没有经典导数的连续函数；下一小节的例子将显示这种描述的用处。

\subsection{逆定理与高阶逼近的边界}
\label{b2:subsec:020904}

先回顾\secref{b2:sec:approximation-inverse}中二进尺度求和的关键位置。误差 \(2^{-j\alpha}\) 乘上导数所需的次数因子 \(2^j\)，得到 \(2^{j(1-\alpha)}\)。当 \(0<\alpha<1\) 时，这个几何级数由最后一项控制；在 \(\alpha=1\) 时，每一项同阶，累积项数变成对数。这说明旧证明只给出 \(h\log(1/h)\) 量级。下面的函数表明这项损失确实可能出现。

\begin{proposition}\label{b2:prop:approximation-integer-counterexample}
函数
\[
 f(x)=\sum_{j=1}^{\infty}2^{-j}\sin(2\pi 2^j x),\quad x\in\mathbb T,
\]
是实值连续周期函数，满足 \(\mathcal E_n(f)\leq2/n\)（整数 \(n\geq1\)），但不是 Lipschitz 连续的。它的二阶连续模却满足 \(\omega_2(f,\delta)\leq C\delta\)（\(0<\delta\leq1/2\)），其中 \(C\) 为绝对常数。
\end{proposition}
\begin{proof}
各项一致范数之和为 \(\sum_{j\geq1}2^{-j}=1\)，故级数一致收敛，和函数连续。若 \(2^N\leq n<2^{N+1}\)，取前 \(N\) 项作为次数不超过 \(n\) 的逼近，尾项范数至多 \(2^{-N}\leq2/n\)；\(n=1\) 时取零多项式。

为检查 Lipschitz 性，取 \(h_N=2^{-N-2}\)。当 \(1\leq j\leq N\) 时，\(2\pi2^jh_N\in[0,\pi/2]\)，由 \(\sin u\geq2u/\pi\) 得
\[
 2^{-j}\sin(2\pi2^jh_N)\geq4h_N.
\]
当 \(j=N+1\) 时正弦自变量为 \(\pi\)，随后为 \(\pi\) 的整数倍，所有尾项恰好为零。于是 \(f(0)=0\) 且
\[
 \frac{|f(h_N)-f(0)|}{h_N}\geq4N\longrightarrow\infty.
\]
不存在统一的 Lipschitz 常数。

另一方面，指数函数的二阶差分乘子是 \((e^{-2\pi\mathrm i 2^jh}-1)^2\)，所以每项的二阶差分范数至多
\[
 2^{-j}\min\{4,(2\pi2^j h)^2\}.
\]
对 \(0<|h|\leq1/2\)，选 \(N\geq1\) 使 \(2^{-N-1}<|h|\leq2^{-N}\)。低频之和至多 \(C h^2\sum_{j\leq N}2^j\leq C|h|\)，高频之和至多 \(4\sum_{j>N}2^{-j}\leq C|h|\)。对 \(|h|\leq\delta\) 取上确界即得结论。
\end{proof}

\(\omega_2(f,\delta)=O(\delta)\) 常用于描述整数光滑阶附近比 Lipschitz 更宽的函数类。这个例子已经足够说明：把上节 Hölder 等价中的 \(\alpha<1\) 改成 \(\alpha\leq1\)，并保持右侧为一阶 Hölder 条件，会得到错误命题。若研究更高整数阶，也必须先选合适的高阶差分尺度，而不是继续使用一阶连续模。

\ProseParagraphBreak

还有一个更基础的限制：若区间上的函数满足 \(|f(x)-f(y)|\leq C|x-y|^\alpha\) 且 \(\alpha>1\)，把任意线段分成 \(m\) 等份可得
\[
 |f(x)-f(y)|\leq C|x-y|^\alpha m^{1-\alpha}\longrightarrow0.
\]
它只能是常数。光滑函数的一阶变化通常仍是线性的，所以高阶光滑性不可能用“一阶差小于距离的高次幂”来定义。高阶差分和导数的 Hölder 条件正是为避免这一矛盾而引入的。

\subsection{端点情形}
\label{b2:subsec:approximation-endpoints}

区间上的代数多项式可以写成 \(p(\cos\theta)\)，再使用三角多项式工具；但 \(x=\cos\theta\) 不是处处保持距离尺度的变换。区间内部 \(\sin\theta\) 不为零，角度变化与位置变化成比例；在 \(x=1\) 附近则有 \(1-\cos\theta\sim\theta^2/2\)。因此同一个角度小量，在端点变成二次的位置小量。

\begin{proposition}\label{b2:prop:approximation-algebraic-derivative}
设 \(p\) 为实或复系数代数多项式，次数不超过整数 \(n\geq1\)，范数取在 \([-1,1]\) 上。则
\[
 \sqrt{1-x^2}\,|p'(x)|\leq n\|p\|_\infty\quad(-1<x<1),
 \quad \|p'\|_\infty\leq\frac\pi2 n^2\|p\|_\infty.
\]
此外，Chebyshev 多项式满足 \(\|T_n\|_\infty=1\) 和 \(T_n'(1)=n^2\)，故全区间导数估计中的 \(n^2\) 不能统一换成 \(n\)。第二个不等式在此只说明次数的正确尺度，不主张常数 \(\pi/2\) 最优。
\end{proposition}
\begin{proof}
令 \(v(\theta)=p(\cos\theta)\)。这是角度周期 \(2\pi\)、次数不超过 \(n\) 的三角多项式。把\cref{b2:thm:approximation-bernstein-trigonometric}中的自变量按 \(\theta=2\pi t\) 换元，得到 \(\|v'\|_\infty\leq n\|p\|_\infty\) 及 \(\|v''\|_\infty\leq n^2\|p\|_\infty\)。链式法则给出
\[
 v'(\theta)=-\sin\theta\,p'(\cos\theta),
\]
因而得到第一个估计。

由于 \(v'(0)=v'(\pi)=0\)，微积分基本定理同时给出
\[
 |v'(\theta)|\leq n^2\|p\|_\infty\min\{\theta,\pi-\theta\}
 \quad(0\leq\theta\leq\pi).
\]
使用 \(\sin\theta\geq(2/\pi)\min\{\theta,\pi-\theta\}\)，相除后得到第二个估计，端点由 \(p'\) 连续性补上。最后，\(T_n(\cos\theta)=\cos(n\theta)\) 给出
\[
 T_n'(1)=\lim_{\theta\to0}\frac{\cos(n\theta)-1}{\cos\theta-1}=n^2,
\]
而 \(|\cos(n\theta)|\leq1\) 并在端点取等号，所以其一致范数为 \(1\)。
\end{proof}

这也影响逆定理的表述。令 \(f(x)=\sqrt{1-x^2}\)（\(-1\leq x\leq1\)）。对应的周期函数 \(g(t)=f(\cos(2\pi t))=|\sin(2\pi t)|\) 是 Lipschitz 函数，常数可取 \(2\pi\)。对 \(g\) 使用周期 Jackson 逼近，再取偶部并把余弦换成 Chebyshev 多项式，就像\secref{b2:sec:approximation-jackson}中的代数转移证明一样，得到 \(E_n(f)\leq C/n\)。然而
\[
 \frac{|f(1-h)-f(1)|}{h^\alpha}
 =h^{1/2-\alpha}\sqrt{2-h}\longrightarrow\infty
 \quad(\alpha>1/2).
\]
因此即使固定 \(1/2<\alpha<1\)，代数最佳误差 \(E_n(f)=O(n^{-\alpha})\) 也不能推出闭区间通常距离下的 \(\alpha\)-Hölder 连续性。周期逆定理没有错；错误发生在忽略了余弦变换的端点退化。

\ProseParagraphBreak

在闭子区间 \([-1+\varepsilon,1-\varepsilon]\) 上，\(\sqrt{1-x^2}\) 有正下界，可以恢复与角度相当的局部尺度。若希望在整个闭区间获得精确的高阶正逆理论，则需要与端点相适应的加权连续模：差分步长也要随位置调整，才能同时反映内部与端点附近的变化。


\begin{chaptersummary}
Bernstein 多项式通过离散正权重把连续函数变为可计算的多项式，二阶矩控制一致误差；Chebyshev 交错定理则通过误差的等幅换号判定实区间上的最佳性，并给出唯一性。两者分别解决“怎样构造”和“怎样判断最优”，不能把一个具体算子的误差自动视为最佳误差。

周期问题中，平方 Fejér 核同时保留有限次数、非负性与单位质量，又改善尾部衰减，给出 Jackson 估计。反向推导用 Bernstein 导数估计控制每个频率尺度的变化，得到 \(0<\alpha<1\) 时 \(\mathcal E_n(f)=O(n^{-\alpha})\) 与 \(\alpha\)-Hölder 连续性的等价。高阶差分继续消去低阶变化，但整数光滑阶须另作判断；余弦换元又在区间端点改变距离尺度。这些限制都是结论的一部分，应当与定理一起记住。
\end{chaptersummary}

\BookExercises

% 来源：自拟。目标：以子空间反例区别有限维最佳存在与多项式空间的唯一性。
\begin{exercise}[子空间中的非唯一性]\label{b2:ex:approximation-nonunique}
在 \(C([-1,1];\mathbb R)\) 中取 \(V=\{ax^2:a\in\mathbb R\}\)。求常函数 \(1\) 到 \(V\) 的一致最佳误差及全部最佳逼近。解释这为什么不与 Chebyshev 交错定理及其唯一性推论矛盾。
\end{exercise}
\begin{solution}
所有 \(ax^2\) 都在 \(x=0\) 取零，所以误差至少为 \(1\)。令 \(t=x^2\in[0,1]\)，仿射函数 \(1-at\) 的绝对值在端点达到最大，因而
\[
 \|1-ax^2\|_\infty=\max\{1,|1-a|\}.
\]
最小值为 \(1\)，且恰好 \(a\in[0,2]\) 时取到。\(V\) 虽然有限维，却不是全部次数不超过某个整数的多项式空间。这里的非零基函数 \(x^2\) 在区间中有零点，不能用“一维空间中非零函数无零点”的性质完成交错唯一性论证。有限维紧性给出最优者存在，唯一性还需要空间的额外结构。
\end{solution}

% 来源：自拟。目标：同一目标函数的二次一致逼近与正交逼近比较。
\begin{exercise}[绝对值函数的两种二次逼近]\label{b2:ex:approximation-absolute-value}
求 \(f(x)=|x|\) 在 \([-1,1]\) 上的最佳一致二次逼近，以及对 Lebesgue 测度的最佳 \(L^2\) 二次逼近。分别计算两者的一致误差。
\end{exercise}
\begin{solution}
先取 \(p(x)=x^2+1/8\)。在 \(0\leq x\leq1\) 上，误差为 \(x-x^2-1/8\)，最小值 \(-1/8\) 在 \(0,1\) 取到，最大值 \(1/8\) 在 \(1/2\) 取到。误差是偶函数，于是它在
\[
 -1,-1/2,0,1/2,1
\]
依次取 \(-1/8,1/8,-1/8,1/8,-1/8\)。任取其中连续的四点即满足二次交错判据。因此 \(p\) 是唯一的最佳一致逼近，误差为 \(1/8\)。

对 \(L^2\) 问题，奇偶正交性使最佳逼近的一次项系数为零。写成 \(q(x)=a x^2+c\)，误差与 \(1,x^2\) 正交给出
\[
 2c+\frac23a=1,\quad \frac23c+\frac25a=\frac12.
\]
解得 \(a=15/16,c=3/16\)。它的 \(L^2\) 最佳性由正交投影定理保证。\(0\leq x\leq1\) 上误差 \(x-15x^2/16-3/16\) 的内部最大值在 \(x=8/15\) 取到，为 \(19/240\)；两个端点的值分别为 \(-3/16,-1/8\)。故 \(\|f-q\|_\infty=3/16>1/8\)。减小平方平均误差会允许某些点的误差更大。
\end{solution}

% 来源：自拟。目标：固定次数反复应用正算子，与次数趋于无穷的逼近相区分。
\begin{exercise}[反复应用 Bernstein 算子]\label{b2:ex:approximation-bernstein-iterates}
固定整数 \(n\geq2\)，以 \(B_n^k\) 表示将 \([0,1]\) 上的 Bernstein 算子复合 \(k\) 次。对任意 \(f\in C([0,1];\mathbb R)\)，证明
\[
 B_n^k f\longrightarrow (1-x)f(0)+xf(1)
\]
一致成立。提示：利用 \(u(x)=x(1-x)\) 及 \(B_nu=(1-1/n)u\)。
\end{exercise}
\begin{solution}
令 \(\ell(x)=(1-x)f(0)+xf(1)\)、\(g=f-\ell\)，则 \(g(0)=g(1)=0\)，且 \(B_n\ell=\ell\)。给定 \(\varepsilon>0\)，由端点连续性，在两端各一小段上 \(|g|\leq\varepsilon\)；剩下的紧区间上 \(u\) 有正下界。因而存在有限的 \(C_\varepsilon\geq0\)，使整个区间上
\[
 |g(x)|\leq\varepsilon+C_\varepsilon u(x).
\]
二阶矩计算给出 \(B_n(x^2)=x^2+x(1-x)/n\)，所以提示中的特征关系成立。正性、保常数及逐次应用产生
\[
 |B_n^k g|\leq B_n^k|g|
 \leq\varepsilon+C_\varepsilon(1-1/n)^ku.
\]
先令 \(k\to\infty\)，再令 \(\varepsilon\downarrow0\)，便得一致收敛。固定采样次数反复平均会逐渐抹去内部形状，只留下端点仿射插值；这与 \(n\to\infty\) 时 \(B_nf\to f\) 是两个不同的极限。
\end{solution}

% 来源：自拟。目标：Fourier系数作逼近下界，构造任意幂次的实际最佳误差。
\begin{exercise}[可精确判断量级的逼近率]\label{b2:ex:approximation-lacunary-rate}
证明：对 \(f\in C(\mathbb T;\mathbb C)\) 及 \(|k|>n\)，有 \(\mathcal E_n(f)\geq|\widehat f(k)|\)。再对任意 \(\alpha>0\)，定义
\[
 f_\alpha(x)=\sum_{j=1}^\infty2^{-j\alpha}\cos(2\pi2^jx).
\]
证明它连续，且存在只依赖 \(\alpha\) 的 \(c_\alpha,C_\alpha>0\)，使 \(c_\alpha n^{-\alpha}\leq\mathcal E_n(f_\alpha)\leq C_\alpha n^{-\alpha}\) 对整数 \(n\geq1\) 成立。
\end{exercise}
\begin{solution}
对任意 \(p\in\mathcal T_n\)，高于次数的系数为零，所以
\[
 |\widehat f(k)|=|\widehat{f-p}(k)|
 \leq\int_{\mathbb T}|f-p|\leq\|f-p\|_\infty.
\]
取下确界得到下界工具。定义 \(f_\alpha\) 的级数由收敛的正项几何级数控制，因而一致收敛，且可以逐项计算 Fourier 系数。

取 \(N\geq0\) 使 \(2^N\leq n<2^{N+1}\)。前 \(N\) 项是可容许的逼近，尾和至多
\[
 \frac{2^{-(N+1)\alpha}}{1-2^{-\alpha}}\leq C_\alpha n^{-\alpha}.
\]
另一方面，在频率 \(2^{N+1}>n\) 处，系数恰好为 \(2^{-1-(N+1)\alpha}\)。由于 \(n\geq2^N\)，它至少为 \(2^{-1-\alpha}n^{-\alpha}\)，给出下界。\(N=0\) 时前零项理解为零函数，论证也成立。这里下界检查了任何次数不超过 \(n\) 的候选者，而不只检查某一种截断算法。
\end{solution}

% 来源：自拟。目标：从Lagrange插值构造最佳误差的有限节点证书。
\begin{exercise}[节点值给出的误差下界]\label{b2:ex:approximation-dual-certificate}
设 \(a\leq x_0<\cdots<x_{n+1}\leq b\)、\(a<b\)，其中 \(n\geq0\) 为整数。令
\[
 w_j=\frac1{\prod_{\ell\ne j}|x_j-x_\ell|},\quad
 \lambda_j=\frac{w_j}{\sum_{\ell=0}^{n+1}w_\ell},\quad
 L(f)=\sum_{j=0}^{n+1}(-1)^{n+1-j}\lambda_j f(x_j).
\]
证明对任意 \(f\in C([a,b];\mathbb R)\)，有 \(E_n(f)\geq|L(f)|\)。说明当这些点是某个最佳误差的交错极值点时，为什么下界取等。
\end{exercise}
\begin{solution}
对 \(p\in\mathcal P_n\)，用 \(n+2\) 个节点写 Lagrange 插值：
\[
 p(x)=\sum_{j=0}^{n+1}p(x_j)
 \frac{\prod_{\ell\ne j}(x-x_\ell)}{\prod_{\ell\ne j}(x_j-x_\ell)}.
\]
右侧与左侧之差次数至多为 \(n+1\)，又在全部 \(n+2\) 个节点为零，故恒为零。比较 \(x^{n+1}\) 的系数得到
\[
 \sum_{j=0}^{n+1}\frac{p(x_j)}{\prod_{\ell\ne j}(x_j-x_\ell)}=0.
\]
分母的符号为 \((-1)^{n+1-j}\)，所以 \(L(p)=0\)。又因 \(\lambda_j>0\)、\(\sum_j\lambda_j=1\)，
\[
 |L(f)|=|L(f-p)|\leq\sum_j\lambda_j|f(x_j)-p(x_j)|
 \leq\|f-p\|_\infty.
\]
取下确界即可。若 \(f(x_j)-p_*(x_j)=\varepsilon(-1)^jE_n(f)\)，则所有带符号的求和项都具有同一符号，得 \(L(f)=\varepsilon(-1)^{n+1}E_n(f)\)。因此有限个取样值便能证明任何候选都不可能优于该误差。
\end{solution}

% 来源：自拟。目标：最佳一致误差经Parseval与分块Cauchy--Schwarz控制Fourier绝对收敛。
\begin{exercise}[逼近率与 Fourier 绝对收敛]\label{b2:ex:approximation-fourier-absolute}
设 \(f\in C(\mathbb T;\mathbb C)\)，且
\[
 \sum_{j=0}^\infty 2^{j/2}\mathcal E_{2^j}(f)<\infty.
\]
证明 \(\sum_{k\in\mathbb Z}|\widehat f(k)|<\infty\)，并证明 Fourier 级数一致收敛到 \(f\)。由此推出：若 \(f\) 是 \(\alpha\)-Hölder 连续函数，\(1/2<\alpha<1\)，则其 Fourier 级数绝对且一致收敛。
\end{exercise}
\begin{solution}
选 \(p_j\in\mathcal T_{2^j}\) 为一致最佳逼近。对块 \(A_j=\{k:2^j<|k|\leq2^{j+1}\}\)，\(p_j\) 的相应系数为零。Bessel 不等式与圆周测度为一给出
\[
 \sum_{k\in A_j}|\widehat f(k)|^2
 \leq\|f-p_j\|_2^2\leq\mathcal E_{2^j}(f)^2.
\]
该块有 \(2^{j+1}\) 个整数，因此 Cauchy--Schwarz 给出
\[
 \sum_{k\in A_j}|\widehat f(k)|\leq2^{(j+1)/2}\mathcal E_{2^j}(f).
\]
对块求和，再加上三个低频 \(k=0,\pm1\)，即得绝对收敛。所得级数一致收敛到某个连续函数 \(g\)，逐项积分说明其 Fourier 系数与 \(f\) 相同。\cref{b2:cor:02-uniqueness}给出几乎处处相等，再由连续性得到处处相等。

若 \(f\) 是所述 Hölder 函数，周期 Jackson 定理给出 \(\mathcal E_{2^j}(f)\leq C2^{-j\alpha}\)，所需级数由 \(\sum_j2^{-j(\alpha-1/2)}\) 控制。该论证在 \(\alpha=1/2\) 时不再给出收敛，因而不能把本题的严格不等号直接删去。
\end{solution}

% 来源：自拟。目标：可求和的导数增量恢复经典可微性，不依赖形式逐项微分。
\begin{optionalexercise}[由更快逼近恢复导数]\label{b2:ex:approximation-derivative-recovery}
设 \(f\in C(\mathbb T;\mathbb C)\)，且 \(\sum_{j=0}^\infty2^j\mathcal E_{2^j}(f)<\infty\)。令 \(p_j\) 为次数不超过 \(2^j\) 的一致最佳逼近。证明 \(f\in C^1(\mathbb T;\mathbb C)\)，且 \(p_j'\to f'\) 一致成立。若进一步 \(\mathcal E_n(f)\leq A n^{-\alpha}\)、\(1<\alpha<2\)，证明 \(f'\) 为 \((\alpha-1)\)-Hölder 连续。
\end{optionalexercise}
\begin{solution}
最佳误差随次数下降，故
\[
 \|p_{j+1}-p_j\|_\infty\leq2\mathcal E_{2^j}(f).
\]
Bernstein 导数估计给出 \(\|p_{j+1}'-p_j'\|_\infty\leq C2^j\mathcal E_{2^j}(f)\)。假设使这些导数增量可求和，所以 \(p_j'\) 一致趋于一个连续周期函数 \(g\)。另一方面，Fejér 一致逼近说明 \(\mathcal E_n(f)\to0\)，从而 \(p_j\to f\) 一致。对 \(x\in[0,1]\)，在
\[
 p_j(x)-p_j(0)=\int_0^xp_j'(t)\,\mathrm dt
\]
中取极限，得到 \(f(x)-f(0)=\int_0^xg(t)\,\mathrm dt\)。因此 \(f'=g\)，端点也按周期延拓相容。这个积分论证才使极限可以求导。

在幂次假设下，对导数增量尾和估计得
\[
 \|f'-p_J'\|_\infty\leq C_\alpha A2^{-J(\alpha-1)}.
\]
若 \(2^J\leq n<2^{J+1}\)，则 \(p_J'\in\mathcal T_n\)，故 \(\mathcal E_n(f')\leq C_\alpha A n^{-(\alpha-1)}\)。由于 \(0<\alpha-1<1\)，对连续函数 \(f'\) 使用主线的周期逆定理，得到所需 Hölder 性。这里先建立导数存在，随后才把逆定理应用于导数。
\end{solution}

% 来源：自拟。目标：把端点导数n²与可见的极值点距离和缩放极限联系起来。
\begin{optionalexercise}[端点附近的振荡尺度]\label{b2:ex:approximation-endpoint-scale}
对 Chebyshev 多项式 \(T_n\)，证明从 \(x=1\) 到相邻极值点的距离为 \(1-\cos(\pi/n)\sim\pi^2/(2n^2)\)。再证明：对任意固定 \(S>0\)，当 \(n\) 足够大时，
\[
 T_n(1-s/n^2)\longrightarrow\cos\sqrt{2s}
\]
在 \(0\leq s\leq S\) 上一致成立，并解释为何端点的自然尺度是 \(n^{-2}\)。
\end{optionalexercise}
\begin{solution}
恒等式 \(T_n(\cos\theta)=\cos(n\theta)\) 给出相邻极值点 \(1\) 和 \(\cos(\pi/n)\)，两点函数值从 \(1\) 变到 \(-1\)。利用 \(1-\cos u=2\sin^2(u/2)\) 即得距离渐近式。

在 \(0\leq s\leq S\) 且 \(n\) 足够大时，令 \(u=\sqrt{s/2}/n\leq1/2\)。由半角公式，
\[
 n\arccos(1-s/n^2)=2n\arcsin u.
\]
由 \(\arcsin u-u=\int_0^u((1-t^2)^{-1/2}-1)\,\mathrm dt\)，且括号在 \([0,1/2]\) 上至多为 \(Ct^2\)，可得
\[
 \bigl|2n\arcsin u-\sqrt{2s}\bigr|\leq C S^{3/2}/n^2.
\]
余弦函数是 Lipschitz 的，于是得到一致收敛。次数虽只有 \(n\)，靠近端点时却能在长为 \(n^{-2}\) 的区间内产生固定大小的变化；这与 \(T_n'(1)=n^2\) 一致，也解释了为何不能把内部的 \(n^{-1}\) 尺度直接移到端点。
\end{solution}
